% P7: Cache Geometry Under Obfuscation
%
% Companion to:
%   P2 Oracle Loop (detection + steering method)
%   P6 Oracle Formulary (therapeutic profiles)
%
% Version 1.1 (October 2026) withdraws version 1.0's claim that KV-Cloak
% blinds external geometric detection and renders adversarial steering
% incoherent. See the correction note after the abstract and the
% Correction Record appendix. Every number in this version is recomputed
% by check_correction_claims.py from the files in data/.

\documentclass[11pt,a4paper]{article}
\usepackage{amsmath,amssymb,booktabs,graphicx,hyperref}
\usepackage{natbib}

\title{Cache Geometry Under Obfuscation:\\
Per-Head Rotation Does Not Hide SVD Features,\\
and a KV-Cloak Implementation Is Inconclusive}

\author{
Thomas Edrington\textsuperscript{1}\\[4pt]
\textsuperscript{1}Liberation Labs\\
Aided by AI research agents.
}

\date{Version 1.1, October 2026}

\begin{document}
\maketitle

% ===================================================================
\begin{abstract}
% ===================================================================

Per-head orthogonal rotation of the key-value (KV) cache, the simulation
of KV-Cloak's obfuscation~\citep{luo2025shadow} used throughout this
paper, does not hide the singular-value features that cache-geometry
detectors read. It leaves the spectrum of each attention head, and of all
heads concatenated along the feature axis, exactly unchanged. In a
100-prompt replication on a Claude-distilled Qwen3.5-27B, all 24
cross-head spectral features and all 16 per-head stable ranks were
unchanged in every trial, and a classifier separating 50
confabulation-inducing prompts from 50 factual questions scored AUROC
0.994 with the rotation, without it, and when fitted without it and
tested with it. Our implementation of the full KV-Cloak mechanism,
$S \cdot \hat{P} \cdot (K + A) \cdot M$, does change the spectrum. On 7
confabulations against 87 hedged replies it reduced the best uncloaked
feature, the spectral gap $\sigma_1/\sigma_2$, from AUROC 0.903 to 0.614,
while signal rank (0.766) and top singular-value excess (0.736) were no
lower than without it. Adjusted for choosing the best of five features,
the strongest remaining feature has $p = 0.054$: at this size the data
show neither that detection survives the mechanism nor that it is
removed. Detection without the cloak is exploratory: one feature on seven
confabulations, and 0.590--0.784 on three smaller models, none significant
after correction. Injection through the cloak was not tested in any
interpretable way, because direct steering, the positive control,
corrected 0--1 of 10 replies per model.

\end{abstract}

\noindent\fbox{\parbox{\dimexpr\linewidth-2\fboxsep-2\fboxrule}{\small
\textbf{Correction note (version 1.1, October 2026).} Version~1.0, titled
``Cache Geometry Under Obfuscation: KV-Cloak as Defense Against Adversarial
KV-Cache Steering'' (in the lab's published-research repository on GitHub
since 29 April 2026; its academic edition has been on Zenodo since June
2026, all versions 10.5281/zenodo.20596922), claimed that KV-Cloak defends
cache-geometry methods against adversaries. Its abstract opened: ``A single
invertible linear transform of the KV cache degrades SVD-based
cognitive-state detection from near-ceiling to near-chance ($p = 5.5 \times
10^{-16}$), while a real KV-Cloak implementation confirms the defense is
conservative (detection AUROC drops from 0.903 to 0.614).'' We withdraw
that sentence and the claims built on it. The $p$-value tests how far a
rotation shifts one feature's values, not detection. The replication
shipped with version~1.0, but not reported in it, found no change at all:
per-head rotation cannot move per-head spectra, and its features did not
move. Under the full mechanism only the spectral gap fell; two other
features stayed at or above their uncloaked values. Also withdrawn:
``obfuscation renders geometric detection and vector injection incoherent
to external adversaries''; ``Obfuscation defeats transfer attacks'' and ``A
detector calibrated on uncloaked caches cannot operate on cloaked caches'',
true only of a fixed threshold in one feature pipeline; that the rotation
result ``represents transfer-attack failure (domain shift under
coordinate change), not signal degradation'', when it was measured with both classes under the same
cloak; that this setting ``is relevant only to an adversary who controls
the cloak key'', when it is the setting of an adversary with no key who
fits a detector on cloaked caches; the heading ``Adversarial Injection Is
Incoherent'', an untested argument presented as a result; that ``KV-Cloak
alone prevents external adversaries from deploying malicious geometric
steering''; and that ``the architecture ensures that defenders retain
capabilities that attackers lose''. Its title asserted a working defense,
so this version has a new one. Corrected numbers: the hedged class is 87
replies, not 89 (six ambiguous replies are excluded, and the Wilcoxon test
reproduces only at 87); the spectral gap's interval is [0.810, 0.980], from
the run file, not [0.806, 0.977], which no file holds; three Table~2 cells
were rounded down (0.754, 0.786 and 0.543 are 0.755, 0.787 and 0.544); the
random-direction results were called ``inconsistent with destabilization
rates above 30\% (power 0.97 to detect)'', but 2 adverse events in 10 allow
rates up to 51\%, and 0.97 is the chance of seeing at least one event; the
positive control's ``3 calibration pairs'' has no record, and the
positive-control file shipped with version~1.0 has 30 pairs and no outcome
labels, so its results are not reported; and the spectral gap, a raw ratio
$\sigma_1/\sigma_2$, was called Marchenko--Pastur-corrected in the
abstract and Results. Version~1.0's text-only baseline (0.755, with the
cache ahead in 96.7\% of bootstrap samples) cannot be traced to any run
file or recomputed from the stored replies, so this version drops it, and
with it any claim about what the cache adds to surface text. The run files
behind both tables, never in this paper's folder, are added. The lab's
remediation plan of 23 June 2026 recommended withdrawing the rotation
claim; it was not applied. The defects were raised again by the lab's
style gate (4 September 2026), by Kavi's verification flags (14 September
2026) and by a scoping review (20 October 2026). The withdrawn claims are
the authors' responsibility. Appendix~\ref{app:record} keeps every
withdrawn passage verbatim.}}

% ===================================================================
\section{Introduction}
% ===================================================================

The discovery that KV-cache geometry encodes cognitive state information
creates a tension between transparency and safety. Publishing detection
methods enables defensive applications (confabulation monitoring,
deception detection, self-regulating AI systems) but simultaneously
provides a roadmap for adversarial exploitation. An attacker with cache
access could:

\begin{itemize}
\item \textbf{Surveil} cognitive states to identify when a model is
uncertain, compliant, or vulnerable to manipulation.
\item \textbf{Inject} emotion vectors calibrated to specific
misalignment types, deliberately inducing confabulation, sycophancy, or
overconfidence.
\item \textbf{Identify} individual users or personas from cache geometry
fingerprints.
\end{itemize}

This paper investigates whether KV-Cloak~\citep{luo2025shadow}---a
lightweight, reversible cache obfuscation scheme designed for privacy
protection---also mitigates these adversarial capabilities. In the forms
we tested, it does not hide the signal. Per-head rotation, the simulation
used here, leaves the spectra these detectors read unchanged
(Section~\ref{sec:replication}). Our implementation of the full mechanism
removes the best detection feature but leaves others at their uncloaked
levels, and seven confabulations are too few to say whether a detectable
signal survives it (Section~\ref{sec:real}).

% ===================================================================
\section{Background}
% ===================================================================

\subsection{The Lyra Technique}

KV-cache geometry---specifically, Marchenko-Pastur-corrected spectral
features derived from SVD of the key tensor---distinguishes cognitive
states during autoregressive generation. Prior work reports
confabulation detection at 0.969--0.999 and hardware invariance
($r>0.999$ across RTX~3090 and H200)~\citep{lyra2026oracle}. A
within-model deception result previously reported at AUROC~1.000 across
seven model scales has since been withdrawn by its authors~\citep{lyra2026techniqueii}: a
same-prompt control collapses it to 0.160, indicating the classifier
separated the system-prompt template rather than deception geometry.
Whether a confound-free within-model deception detector is achievable
remains open; nothing in this paper depends on it. Contrastive activation addition
(CAA)~\citep{panickssery2023caa} enables injection of emotion vectors into the cache to steer model
behavior, correcting confabulation at rates up to 96\% (hostile vector
on Qwen2.5-7B-Instruct, $n=68$, McNemar $p<0.0001$)~\citep{lyra2026formulary}.

\subsection{KV-Cloak}

Luo et al.\ demonstrate that an attacker can reconstruct sensitive user
inputs directly from the KV cache via inversion, collision, and
injection attacks. KV-Cloak mitigates this with a reversible
matrix-based obfuscation scheme combined with operator fusion, reducing
reconstruction quality to random noise with minimal impact on model
accuracy. The obfuscation operates inside a trusted execution
environment (TEE); external consumers see only the obfuscated cache.

\subsection{Threat Model}

We consider an adversary who has read and/or write access to the
KV cache during inference (e.g., through a shared serving infrastructure,
a compromised inference endpoint, or a supply-chain attack on the
serving stack). The adversary's goal is to deploy a malicious Oracle
Loop: detect vulnerable cognitive states via geometric features, then
inject steering vectors to destabilize the model.

The defender may optionally run a legitimate Oracle Loop inside the TEE
for misalignment monitoring. The question this paper tests, whether
obfuscation blinds external geometric methods, does not depend on whether
a defensive Oracle is deployed.

% ===================================================================
\section{Method}
% ===================================================================

\subsection{Obfuscation Simulation}
\label{sec:method_rot}

We simulate KV-Cloak's reversible matrix obfuscation as per-head
orthogonal rotation: for each attention head independently, we generate
a random orthogonal matrix $Q_h$ and apply $K'_h = K_h Q_h$,
$V'_h = V_h Q_h$. This leaves the singular values of each head's key
matrix unchanged, and also those of the heads concatenated along the
feature axis, because the block-diagonal matrix of the $Q_h$ is itself
orthogonal. It can change the spectrum only of a matrix that combines the
heads in another way, for example stacked along the token axis.

We additionally test:
\begin{itemize}
\item \textbf{Global orthogonal rotation}: a single $Q$ shared across
all heads per layer (preserves cross-head structure); run on the base
model only.
\item \textbf{Random invertible (non-orthogonal)}: $M = I + \sigma N$
where $N$ is Gaussian.
\item \textbf{Additive noise}: Gaussian noise at 1\%, 5\%, 10\%, 50\%
of cache norm.
\item \textbf{Diagonal scaling}: per-dimension random scaling.
\item \textbf{No obfuscation}: baseline control.
\item \textbf{Full KV-Cloak}: our implementation of the whole mechanism
(Section~\ref{sec:real}).
\end{itemize}

\subsection{Feature Extraction}

For each trial, we extract five Marchenko-Pastur-derived features from
the full KV cache (encoding + generation phases): signal rank, signal
fraction, top singular value excess, spectral gap, and norm per token.
The spectral gap feature ($\sigma_1/\sigma_2$ ratio) is not itself
MP-corrected; it is a raw singular value ratio included alongside the
MP-corrected features. The evaluation applies to these feature pipelines
and to confabulation detection only; deception and other cognitive-state
classifiers were not tested under obfuscation.
Features are FWL-residualized \citep{frisch1933partial,lovell1963seasonal} against output token count and evaluated
via Mann-Whitney AUROC with 2000-iteration bootstrap confidence
intervals. AUROCs are direction-free, $\max(A, 1-A)$, so chance gives
values above 0.5 at small $n$; Section~\ref{sec:real} reports a
permutation test that accounts for this and for picking the best of five
features. The script that extracted these features has not been
recovered; the per-trial features and results it wrote are in this
paper's \texttt{data/} folder (Data, Code and Review History).

\subsection{Replication}
\label{sec:method_rep}

A separate replication (\texttt{code/kv\_cloak\_replication.py},
1,038 lines) ran 100 prompts through the distilled model: 50 about
fabricated entities or future events and 50 factual questions. At the
first four full-attention layers (3, 7, 11, 15) it computed six SVD
features (stable rank, spectral entropy, top singular-value ratio, norm,
$\sigma_1$, condition number) from the four heads' keys concatenated along
the feature axis, and each head's stable rank, before and after one set of
four Haar-random head rotations (seed 2026) used for every trial. A
logistic-regression classifier with grouped five-fold cross-validation was
fitted to separate the two prompt types.

\subsection{Models and Prompts}

\begin{itemize}
\item \textbf{Distilled}: Qwen3.5-27B-Claude-4.6-Opus-Reasoning-Distilled
(Jackrong). Hybrid architecture, 16 full-attention + 48 GatedDeltaNet
layers. 100 confab-inducing prompts; behavioral classification via the
Oracle Formulary classifier (40+ hedge patterns). Yields $n=7$
confabulations and $n=87$ hedged replies; 6 ambiguous replies are
excluded.

\item \textbf{Base}: Qwen2.5-7B-Instruct. Standard attention, 28
layers. Same prompt set; yields $n=48$ confabulations, $n=52$ hedged.
\end{itemize}

% ===================================================================
\section{Results}
% ===================================================================

\subsection{Detection on Uncloaked Cache}

On the distilled model, the spectral gap ($\sigma_1/\sigma_2$) achieves
AUROC~0.903 (95\% bootstrap CI: [0.810, 0.980]; FWL-corrected: 0.877,
$R^2 = 0.068$ vs token count; $n=7$ confabulations, $n=87$ hedged)
for confabulation detection using full-cache features. It is the best of
five features and the only one that survives correction for five
($p < 0.01$); adjusted for that choice by label permutation, $p < 0.001$.
Signal rank also detects (0.755 [0.578, 0.918]); the other three features
are not significant ($p \geq 0.06$).

On the base model ($n=48$ confabulations, $n=52$ hedged), the best
feature (signal fraction) achieves only 0.661 (FWL: 0.687,
$R^2 = 0.602$). Detection is weak and heavily confounded with token
count. This difference between models is consistent with multiple
explanations including model size (27B vs 7B), architecture
(hybrid vs standard), and training (Claude-distilled vs
instruction-tuned). We cannot isolate the contribution of RLHF
calibration from these confounds.

\subsection{Rotation Leaves the Replication's Features Unchanged}
\label{sec:replication}

In the replication (Section~\ref{sec:method_rep}), all 24 cross-head
features (six features at four layers) and all 16 per-head stable ranks
were unchanged by the rotation in every one of the 100 trials, to
floating-point precision (largest relative change below $10^{-6}$). The
classifier scored AUROC 0.994 on uncloaked features and 0.994 on cloaked
features (the run file records 0.9944 for both, and a drop of 0.0), and
0.994 when fitted on uncloaked caches and tested on cloaked ones. Its
labels are prompt types, not behaviour: among the 50
confabulation-inducing prompts the behaviour field reads 35 confabulated
and 15 hedged. So this is not a confabulation detector. But whatever
these features detect, they detect identically under the rotation, with
or without retraining. A rotation-based cloak therefore does not hide them
from an adversary who computes them per head or across heads in this
layout.

\subsection{The Rotation Run Behind Table~2: A Shift in Scale, Not a Loss of Ranking}
\label{sec:shift}

In the run behind Table~\ref{tab:real_vs_sim} (\texttt{data/kv\_cloak\_v2}),
the same kind of rotation did change the features. The hedged replies'
spectral gap fell from $24.7 \pm 3.7$ to $10.8 \pm 1.3$ (Wilcoxon
signed-rank $W = 0$, $p = 5.5 \times 10^{-16}$, $n = 87$ paired trials);
every reply moved down. The run's feature script was not recovered, so we
cannot say how it arranged the heads; by the argument in
Section~\ref{sec:method_rot}, it cannot have been per head or along the
feature axis. The rotation result, 0.787, reproduces from this run's
per-trial features: the difference from the replication lies in how the
features were computed, not in the rotation.

The shift breaks a fixed threshold. The spectral-gap threshold that best
separates the classes without the cloak (Youden's $J$), which flags all 7
confabulations and 21 of the 87 hedged replies, flags none of the 94 replies under the rotation, and all
94 under the full mechanism. It does not break the ranking: within the
rotation, the spectral gap still ranks confabulations above hedged
replies, in the same direction, at 0.787. An adversary who refits a
threshold on cloaked caches is not stopped by a change of scale, and one
who computes per-head features sees no change at all
(Section~\ref{sec:replication}).

\subsection{Adversarial Injection Was Not Tested}

CAA~\citep{panickssery2023caa} steering vectors are directions in cache space extracted via
contrastive activation: $d = \mathbb{E}[\text{emotional}] -
\mathbb{E}[\text{neutral}]$. Injection adds $\alpha \cdot d$ to the
cache keys and values. Under per-head rotation $Q_h$, a vector $d$ added
to the cloaked cache becomes $d \cdot Q_h^{-1}$ in the model's own basis,
a different direction if $Q_h$ is unknown to the attacker. That is an
argument, not a result. The test of it was uninterpretable: direct
steering, the positive control, corrected 0--1 of 10 replies on each of
four models, and on Mistral a random direction of the same norm gave
exactly the same counts as the steering vector (1 corrected, 2 adverse,
7 unchanged). An injection that does nothing without the cloak cannot show
what the cloak does to it.

\subsection{A Defender Inside the TEE}

A defender operating inside the TEE has access to the pre-obfuscation
cache, so its detection is the uncloaked detection (AUROC~0.903 on the
distilled model), and its steering vectors are injected in the model's own
basis. The obfuscation is applied after Oracle processing, so the
defensive loop is unaffected. This holds by construction and was not run
as a system. It gives defenders an advantage only if attackers lose the
signal, which the results above do not show.

\subsection{Within-Cloak Detection}
\label{sec:withincloak}

With both classes under the same rotation, detection on the distilled
model is 0.787 (Table~\ref{tab:real_vs_sim}). This is the setting of an
adversary who has no key but can label some cloaked caches, for example
from the model's outputs, and fit a detector on them; an adversary who
holds the key can remove the cloak and use the uncloaked 0.903. Across the
five models with a detection result (Table~\ref{tab:cross_arch}),
within-rotation detection fell on four and rose on one: on Qwen2.5-14B it
went from 0.784 to 0.951 ($p < 0.001$). We have no explanation for the
rise. On Mistral it
fell from 0.707 to 0.516, but the uncloaked 0.707 does not survive
correction for four comparisons. Under each of the other transforms run
on the distilled model (random invertible, diagonal scaling, Gaussian
noise at four levels), the best of the five features stayed at 0.765 or
above.

\subsection{Cross-Architecture Results}

Detection and obfuscation effects vary across architectures
(Table~\ref{tab:cross_arch}).

\begin{table}[h]
\centering
\caption{Cross-architecture KV-Cloak results. Prompts per model: 100 for Qwen3.5-27B-Dist.\ ($n=7$ confabulations) and Qwen2.5-7B, 50 for the other four. Det.(clk) is detection with both classes under the per-head rotation.}
\label{tab:cross_arch}
\begin{tabular}{lrrrr}
\toprule
Model & Confab\% & Det. & Det.(clk) & Mag. adv. \\
\midrule
Qwen3.5-27B-Dist. & 7\% & 0.903 & 0.787 & --- \\
Qwen2.5-14B & 12\% & 0.784 & 0.951$^*$ & 1/10 \\
Mistral-7B & 26\% & 0.707 & 0.516 & 2/10 \\
Llama-3.1-8B & 14\% & 0.590 & 0.559 & 1/10 \\
Qwen2.5-7B & 48\% & 0.661 & 0.639 & --- \\
DeepSeek-R1-7B & 2\% & --- & --- & 0/10 \\
\bottomrule
\multicolumn{5}{l}{\footnotesize $^*$Rises under the rotation; unexplained.} \\
\multicolumn{5}{l}{\footnotesize Mag.\ adv.\ = adverse events from magnitude attack (random direction, 10 trials).}
\end{tabular}
\end{table}

Uncloaked detection ranges from 0.590 (Llama) to 0.903 (distilled Qwen);
no cross-architecture model survives correction for four comparisons.
Perturbation along a random direction, at the norm of a steering vector,
produced 0--2 adverse events per model in 10 trials, 4 of 40 pooled over the four
models (one-sided 95\% upper bound 21\%). Per model, 10 trials exclude
little: 2 adverse events in 10 are compatible with rates up to 51\%.

\subsection{Simulation vs Real KV-Cloak Implementation}
\label{sec:real}

We implemented the full KV-Cloak mechanism
($K' = S \cdot \hat{P} \cdot (K + A) \cdot M$) following
Luo et al.\ and compared it against our per-head orthogonal rotation
simulation on the same 100 prompts (Table~\ref{tab:real_vs_sim}). Our
implementation follows their description; it is not the authors' code.

\begin{table}[h]
\centering
\caption{Detection AUROC: uncloaked vs simulation vs real KV-Cloak
(distilled model, $n=7$ confab, $n=87$ hedged). Bootstrap 95\% CIs.}
\label{tab:real_vs_sim}
\begin{tabular}{lrrr}
\toprule
Feature & Uncloaked & Sim.\ rotation & Real KV-Cloak \\
\midrule
spectral\_gap & \textbf{0.903} [.81,.98] & 0.787 [.58,.95] & \textbf{0.614} [.51,.81] \\
signal\_rank & 0.755 [.58,.92] & 0.770 [.62,.90] & 0.766 [.57,.94] \\
signal\_fraction & 0.714 [.52,.90] & 0.683 [.51,.88] & 0.681 [.51,.88] \\
top\_sv\_excess & 0.609 [.51,.79] & 0.544 [.50,.76] & 0.736$^*$ [.53,.92] \\
norm\_per\_token & 0.598 [.51,.77] & 0.598 [.51,.77] & 0.673 [.52,.85] \\
\bottomrule
\multicolumn{4}{l}{\footnotesize $^*$Rises under the mechanism; see text.}
\end{tabular}
\end{table}

The full mechanism moves $\sigma_1$: the additive mask raises the spectral
gap of every reply, from a hedged mean of 24.7 to 84.6 or more. The
spectral gap, the best feature without the cloak, falls to 0.614
($p = 0.33$). Signal rank (0.766, $p = 0.020$) and top singular-value
excess (0.736, $p = 0.038$) are no lower than without the cloak (0.755 and
0.609). No feature survives correction for five ($p < 0.01$), and the best
of the five, 0.766, is about what selection alone produces at this size:
with seven confabulations, the best of five features reaches 0.77 in 5\%
of label shuffles, so its selection-adjusted $p = 0.054$ (under the
rotation, 0.787 with $p = 0.042$; without a cloak, 0.903 with
$p < 0.001$). The mechanism removes the one feature that detected well.
Whether the cache still carries a detectable confabulation signal under
it, these data cannot say.

\paragraph{Consistent treatment of the $\sigma_1$-dependent features.}
We flag that spectral\_gap ($\sigma_1/\sigma_2$) and top\_sv\_excess are
both $\sigma_1$-dependent, and the additive mask $A$ perturbs $\sigma_1$
in both---so the same mechanism drives the spectral\_gap \emph{decrease}
(0.903$\to$0.614) and the top\_sv\_excess \emph{increase} (0.609$\to$0.736).
We should not read one as ``real defense'' and the other as ``artifact''
on different terms. The honest reading is unified: the mask perturbs
$\sigma_1$, which disrupts $\sigma_1$-dependent detection
features---degrading the feature that happened to detect well
(spectral\_gap) and reshuffling one that did not (top\_sv\_excess). Both
movements are mask-on-$\sigma_1$ effects; neither should be cited as a
topology-preserving guarantee, and the spectral\_gap degradation in
particular should be read as ``the mask disrupts this feature,'' not as
evidence the underlying confabulation geometry was destroyed.

% ===================================================================
\section{Discussion}
% ===================================================================

\subsection{What This Means for Dual Use}

Version~1.0 argued that KV-Cloak mitigates the Lyra Technique's dual-use
risk, so that cache-geometry methods could be published while deployed
models stayed protected. The results do not support that. A
rotation-based cloak leaves per-head spectra unchanged, so it offers no
protection against an adversary who computes them, and a shift in scale
in one feature pipeline stops only a detector whose threshold is never
refitted. Our implementation of the full mechanism removed the best
feature, but with seven confabulations we cannot tell whether it removed
the signal. Nothing shown here mitigates the dual-use risk of publishing
cache-geometry methods. A defense would have to be shown on its own terms:
against adaptive adversaries who choose their features after seeing the
cloak, with enough confabulations to measure what remains.

\subsection{What a Test of the Architecture Would Need}

The deployment architecture version~1.0 recommended remains a proposal,
untested here:

\begin{enumerate}
\item KV-Cloak obfuscation active on all external cache access.
\item Oracle Loop (if deployed) operating inside the TEE on the
pre-obfuscation cache.
\item Obfuscation key material and Oracle parameters both
TEE-protected.
\item Per-session or per-query key rotation to prevent statistical
recovery from accumulated observations.
\end{enumerate}

The replication used one rotation for every trial, so per-query key
rotation (item 4) was never exercised, and nothing here measured what it
costs an adversary to recover the signal.

\subsection{Limitations}

\begin{enumerate}
\item \textbf{Small sample size.} The distilled model yields only
$n=7$ confabulations from 100 prompts. The headline AUROC 0.903 has
bootstrap CI [0.810, 0.980] and is the best of five features; adjusted for
that choice it remains significant, but the comparisons under the cloak
(Section~\ref{sec:real}) are not powered. Confirmation on a
higher-confab-rate model is needed.

\item \textbf{Feature code not recovered.} The script behind
Tables~\ref{tab:cross_arch} and~\ref{tab:real_vs_sim} has not been
recovered. Its per-trial outputs are in \texttt{data/}, but how
it arranged heads and layers, which decides whether a rotation can move
its features, cannot be checked.

\item \textbf{Cache versus text untested.} Version~1.0's text-only
baseline cannot be traced or recomputed (Correction note), so nothing here
shows what the cache adds to surface text features.

\item \textbf{TEE trust assumption.} The defensive architecture assumes
TEE integrity. Documented side-channel attacks against SGX
(Van~Bulck et al., 2018; Schwarz et al., 2019) and SEV
(Li et al., 2021) mean this is a hardware trust assumption, not a
cryptographic guarantee. The security reduction is to TEE integrity,
not to a standard hardness assumption.

\item \textbf{Injection test underpowered.} Direct steering, the
positive control, corrected only 0--1 of 10 replies per model, too few to
establish a working baseline. Through-cloak injection results are
therefore uninterpretable: we cannot distinguish ``cloak defeats
injection'' from ``injection was never effective.'' The 30-pair
positive-control run shipped with this paper
(\nolinkurl{data/kv_cloak_positive_control.json}) stored replies without
outcome labels and is not analysed here.

\item \textbf{Cross-architecture significance.} No individual model in
the cross-architecture sweep achieves significance after Bonferroni
correction for four comparisons. Detection is individually suggestive on
Qwen-14B (0.784) and Mistral (0.707) but should not be interpreted as
confirmed cross-architecture replication.

\item \textbf{Adaptive adversary untested.} All tests assume a static,
non-learning attacker. An adversary who can submit probe queries and
observe behavioral outputs could iteratively estimate the rotation
basis. Formal query-complexity bounds are not provided. The threat model
is bounded to single-query, non-adaptive adversaries.
\end{enumerate}

% ===================================================================
\section{Conclusion}
% ===================================================================

Per-head rotation, the simulation of KV-Cloak this paper relied on, does
not hide the spectral features that cache-geometry detectors read: in the
replication none of them moved, and a detector fitted without the rotation
worked unchanged with it. Our implementation of the full mechanism removed
the best detection feature on the distilled model but left two others at
their uncloaked levels, and with seven confabulations we cannot say
whether a detectable signal remains. Version~1.0 concluded that KV-Cloak
and the Oracle Loop are complementary, with defenders keeping what
attackers lose; nothing here shows that attackers lose the signal. Whether
any reversible obfuscation blinds cache-geometry detection is open, and a
test of it needs adaptive adversaries, per-query keys and many more
confabulations.

\section*{Data, Code and Review History}

{\raggedright
Tables~\ref{tab:cross_arch} and~\ref{tab:real_vs_sim}:
\nolinkurl{data/kv_cloak_v2/}, \nolinkurl{data/kv_cloak_real/},
\nolinkurl{data/kv_cloak_cross_arch/} and \nolinkurl{data/kv_cloak_base/},
added in version~1.1, copied unchanged from the lab's KV-Experiments
repository (\nolinkurl{results/oracle_harness/}, commit 8429260, May 2026;
that repository is private). The scripts that wrote them were not
recovered. Replication: \nolinkurl{code/kv_cloak_replication.py} with
\nolinkurl{data/kv_cloak_replication_100.json} and
\nolinkurl{data/kv_cloak_replication_analysis.json}. Positive control:
\nolinkurl{data/kv_cloak_positive_control.json}. Every number in this
version is recomputed from these files by
\nolinkurl{check_correction_claims.py}, in this paper's folder of
\url{https://github.com/Liberation-Labs-THCoalition/published-research}.
Review history: the lab's remediation plan (23 June 2026) recommended
withdrawing the rotation claim; the Agni style gate (4 September 2026)
rejected version~1.0
(\nolinkurl{tools/agni/style/kv-cloak-defense-paper.agni_style.json});
Kavi's verification flags (14 September 2026); for version~1.1, a scoping
review (20 October 2026), the correction gate above, and the Agni paper
and style modes (records in \texttt{review/}).\par}

\bibliographystyle{plainnat}
\section*{Acknowledgments}
We thank Dwayne Wilkes (Sentient Futures / Liberation Labs) for statistical auditing and red-team review.

\bibliography{../references}

% ===================================================================
% Appendix: correction record
% ===================================================================
\appendix
\sloppy

\section{Correction Record (October 2026)}
\label{app:record}

This appendix keeps the withdrawn wording of version~1.0 visible. Each
item quotes it verbatim and says where the evidence is.

\begin{enumerate}
\item \textbf{Title.} ``Cache Geometry Under Obfuscation: KV-Cloak as
Defense Against Adversarial KV-Cache Steering.'' It asserted a working
defense (Sections~\ref{sec:replication} and~\ref{sec:real}).

\item \textbf{Abstract, opening.} ``A single invertible linear transform
of the KV cache degrades SVD-based cognitive-state detection from
near-ceiling to near-chance ($p = 5.5 \times 10^{-16}$), while a real
KV-Cloak implementation confirms the defense is conservative (detection
AUROC drops from 0.903 to 0.614).'' The $p$-value is the Wilcoxon test of
Section~\ref{sec:shift}, a shift in one feature's values. The replication
found no change at all (Section~\ref{sec:replication}), and under the full
mechanism only the spectral gap fell, while signal rank stayed at 0.766
(Section~\ref{sec:real}).

\item \textbf{Abstract, finding 1 (text baseline).} ``A text-only
surface-feature baseline achieves 0.755; cache geometry outperforms text
in 96.7\% of bootstrap samples (delta 0.148, 95\% CI [$-$0.008,
0.322])---suggestive but not significant at the 95\% level.'' Results:
``The best text-only feature (average word length) achieves 0.755 [0.560,
0.924]. Cache geometry outperforms text in 96.7\% of bootstrap samples,
but the 95\% CI on the difference includes zero (delta: 0.148 [$-$0.008,
0.322]).'' No run file holds these numbers, and the stored trials keep
only the first 150 characters of each reply.

\item \textbf{Abstract, finding 2.} ``\textbf{Obfuscation defeats
transfer attacks.} [\ldots] A detector calibrated on uncloaked caches
cannot operate on cloaked caches. This is a domain-shift effect: rotation
preserves the underlying geometric structure but changes the coordinate
system, preventing transfer from uncloaked calibration data.'' True only
of a fixed threshold in one feature pipeline; in the replication's
pipeline nothing moves, and the ranking survives in both
(Sections~\ref{sec:replication} and~\ref{sec:shift}).

\item \textbf{Abstract, finding 3.} ``0--2 adverse events per model in 10
trials, inconsistent with destabilization rates above 30\% (power 0.97 to
detect) but unable to rule out rates up to 20\% at this sample size'' and
``(3 calibration pairs; direct injection produced 0--1/10 corrections)''.
Two events in 10 allow rates up to 51\%; 0.97 is $1 - 0.7^{10}$, the
chance of at least one event; no record gives 3 calibration pairs.

\item \textbf{Abstract, closing.} ``We conclude that KV-Cloak and the
Oracle Loop are complementary: the Oracle Loop provides misalignment
detection and correction; KV-Cloak prevents external adversaries from
weaponizing the same geometric methods. Together, they address the
dual-use risk inherent in publishing cache-geometry techniques.'' And:
``The rotation-simulation result (0.786) represents transfer-attack
failure (domain shift under coordinate change), not signal
degradation.'' The 0.786 (0.787) was measured with both classes under the
same rotation, not by transferring a detector.

\item \textbf{Introduction.} ``We show that it does: obfuscation renders
geometric detection and vector injection incoherent to external
adversaries, while a defender operating inside the trusted execution
environment retains full access to the uncloaked cache and its geometric
properties.''

\item \textbf{Threat Model.} ``[\ldots] but our primary result---that
KV-Cloak blinds external geometric methods---holds regardless of whether
a defensive Oracle is deployed.''

\item \textbf{Method.} ``This preserves within-head singular values while
scrambling cross-head correlations.'' Concatenated along the feature
axis, the heads' singular values are preserved too
(Section~\ref{sec:method_rot}).

\item \textbf{Results, feature-space shift.} ``The cloaking is perfectly
detectable by paired comparison---but this is the desired property. An
external adversary's detection model, calibrated on uncloaked caches,
encounters a completely alien feature distribution. Any detection
thresholds, classifier boundaries, or anomaly scores calibrated on
uncloaked data are meaningless on cloaked data.'' (Section~\ref{sec:shift}.)

\item \textbf{Results, injection.} The heading ``Adversarial Injection Is
Incoherent'' and ``The injection becomes noise rather than a targeted
perturbation.'' Untested: the positive control did not work.

\item \textbf{Results, within-cloak detection.} ``On Mistral, the cloak
effectively eliminates detection entirely. [\ldots] This scenario---where
both classes share the same obfuscation---is relevant only to an
adversary who controls the cloak key, which the TEE architecture
prevents.''

\item \textbf{Results, cross-architecture.} ``Obfuscation reduces
detection on Mistral to chance (0.516) and partially degrades on the
distilled model (0.786 under simulation). Magnitude-based attacks
(direction-agnostic perturbation) produce at most 2 adverse events in 10
trials across all architectures---models are robust to random
perturbation at steering-vector magnitude.''

\item \textbf{Results, simulation vs implementation.} ``Both degrade a
transfer attack---a detector calibrated on uncloaked caches applied to
cloaked caches---rather than destroying the underlying geometric
signal.'' On top singular-value excess: ``This is an artifact of the mask
magnitude, not a genuine detection signal, and would not generalize to
unseen mask instances under per-query key rotation''; the table note
``Anomalous increase; likely artifact of additive mask $A$ inflating top
singular values differentially across behavioral classes''; and ``the
defense works \emph{by} perturbing $\sigma_1$''. Per-query rotation was
not tested.

\item \textbf{Discussion.} ``The Lyra Technique's dual-use risk---that
published cache-geometry methods could enable adversarial steering---is
mitigated by KV-Cloak. Cache obfuscation renders both geometric detection
and vector injection incoherent to external adversaries. A model operator
deploying KV-Cloak can publish Oracle Loop methods transparently while
protecting deployed models from exploitation via those same methods.''
``This addresses the responsible-disclosure tension identified by
Watson~(2026): the defense can be published alongside the technique,
enabling the research community to build on the methods while preventing
weaponization against production systems.'' ``Without a defensive Oracle
Loop, KV-Cloak alone prevents external adversaries from deploying
malicious geometric steering.''

\item \textbf{Conclusion.} ``KV-Cloak and the Oracle Loop are
complementary technologies. The Oracle Loop enables detection and
correction of misalignment through KV-cache geometry. KV-Cloak prevents
external adversaries from exploiting those same geometric methods.
Together, they address the dual-use risk inherent in publishing
cache-geometry research: the technique is transparent, the defense is
available, and the architecture ensures that defenders retain
capabilities that attackers lose.''

\item \textbf{Academic editions only.} The abstract of the academic file
\nolinkurl{academic_main.tex}: ``degrades the strongest detection feature
(spectral gap) from 0.903 to 0.614---near chance---confirming that the
defense is stronger than our initial orthogonal-rotation simulation
(0.786) suggested.'' Its Results: ``confirming that per-head orthogonal
rotation alone \emph{understates} the defense. [\ldots] The simulation was
conservative: the actual defense is stronger than we initially
reported.'' A limitation in both academic files: ``The real mechanism
degrades spectral gap detection more than the simulation (0.614 vs
0.786), confirming that the simulation was conservative.'' Another, in
both: ``Until resolved, the scope of the obfuscation-degrades-detection
claim excludes this model'' (Qwen2.5-14B); the claim it scoped is
withdrawn, and the rise is reported in Section~\ref{sec:withincloak}.
\end{enumerate}

\end{document}
